% 来源: 19c9fbc1-c952-8bb0-8000-000023e32f4d_4.1_体系结构.pdf
% 提取时间: 2026-02-28T22:51:43+08:00

4.1 体系结构
1.   体系结构
我们一直在追求的是在保证底层硬件性能（performance）的同时，使系统易于使用（easy to use）。
性能包含稳定性，可扩展能力，以及快速迭代的能力。易用性来自于使用强大的系统抽象。在实际部署
过程中，性能和可用性往往存在tradeoff，例如过分抽象可能会因为不能做全栈优化带来巨大的性能损
失。
  网络体系结构
1.1.

首先，当电脑第一次联网，由于人们首先想到的是让两台或多台电脑相互通讯，因此，大家构思出
了如下图的逻辑。




互相之间可以通讯的两个服务可以满足最终用户的一些需求，让我们把这个图再详细一点，添加一
些网络协议栈组件：




上述这个修改过的模型自20世纪50年代以来一直使用至今。一开始，计算机很稀少，也很昂贵，
所以两个节点之间的每个环节都被精心制作和维护。随着计算机变得越来越便宜，连接的数量和数
                                                       1
据量大幅增加。随着人们越来越依赖网络系统，工程师们需要保证他们构建的软件能够达到用户所
要求的服务质量。当然，还有许多问题急需解决以达到用户要求的服务质量水平。人们需要找到解
决方案让机器互相发现、通过同一条线路同时处理多个连接、允许机器在非直连的情况下互相通
信、通过网络对数据包进行路由、加密流量等等。
在这其中，有一种叫做流量控制的东西，下面我们以此为例。流量控制是一种防止一台服务器发送
的数据包超过下游服务器可以承受上限的机制。这是必要的，因为在一个联网的系统中，你至少有
两个不同的、独立的计算机，彼此之间互不了解。计算机A以给定的速率向计算机B发送字节，但
不能保证B可以连续地以足够快的速度来处理接收到的字节。例如，B可能正在忙于并行运行其他
任务，或者数据包可能无序到达，并且B可能被阻塞以等待本应该第一个到达的数据包。这意味着
A不仅不知道B的预期性能，而且还可能让事情变得更糟，因为这可能会让B过载，B现在必须对所
有这些传入的数据包进行排队处理。
一段时间以来，大家寄希望于建立网络服务和应用程序的开发者能够通过编写代码来解决上面提出
的挑战。在我们的这个流程控制示例中，应用程序本身必须包含某种逻辑来确保服务不会因为数据
包而过载。在我们的抽象示意图中，它是这样的：




幸运的是，技术的发展日新月异，随着像TCP/IP这样的标准的横空出世，流量控制和许多其他问
题的解决方案被融入进了网络协议栈本身。这意味着这些流量控制代码仍然存在，但已经从应用程
序转移到了操作系统提供的底层网络层中。
这个模型相当地成功。几乎任何一个组织都能够使用商业操作系统附带的TCP/IP协议栈来驱动他
们的业务，满足高性能和高可靠性的要求。
多年以来，计算机变得越来越便宜，并且到处可见，而上面提到的网络协议栈已被证明是用于可靠

                                            2
连接系统的事实上的工具集。随着节点和稳定连接的数量越来越多，行业中出现了各种各样的网络
系统，从细粒度的分布式代理和对象到由较大但重分布式组件组成的面向服务的架构。
历史往往会重演，为了处理更复杂的问题，需要转向更加分散的系统（我们通常所说的微服务架
构），这在可操作性方面提出了新的要求。下面则列出了一个必须要处理的东西：
①计算资源的快速提供
②基本的监控
③快速部署
④易于扩展的存储
⑤可轻松访问边缘
⑥认证与授权
⑦标准化的RPC




因此，尽管数十年前开发的TCP/IP协议栈和通用网络模型仍然是计算机之间相互通讯的有力工
具，但更复杂的架构引入了另一个层面的要求，这再次需要由在这方面工作的工程师来实现。例
如，对于服务发现和断路器，这两种技术已用于解决上面列出的几个弹性和分布式问题。
 ● 服务发现是指在满足特定查询条件的前提下，自动定位服务实例的过程。例如，一个名


                                               3
    为“Teams”的服务可能需要查找一个名为“Players”的服务实例，且该实例
    的 environment属性需设置为 production。此时，通过调用服务发现机制，系统将返
    回一个符合指定条件的服务实例列表。在集中式架构中，这一任务相对简单，通常可以通过诸
    如DNS、负载均衡器以及端口约定（如所有服务默认将HTTP服务器绑定至8080端口）等传统
    手段实现。然而，在分布式或高度分散的环境中，服务发现的复杂性显著增加。过去依赖于
    DNS解析来获取依赖关系的方式，已无法满足现代分布式系统的需求。当前，服务需要处理诸
    如客户端负载均衡、多环境配置、地理位置分布等复杂的挑战。相比于以往仅需一行代码即可
    完成主机名解析的场景，如今的服务实现可能需要引入大量额外逻辑，以应对分布式架构带来
    的各种问题，包括但不限于服务注册、健康检查、动态路由和容错机制等。
 ● 断路器背后的基本思路非常简单。将一个受保护的函数调用包含在用于监视故障的断路器对象
    中。一旦故障达到一定阈值，则断路器跳闸，并且对断路器的所有后续调用都将返回错误，并
    完全不接受对受保护函数的调用。通常，如果断路器发生跳闸，你还需要某种监控警报。
这些都是非常简单的设备，它们能为服务之间的交互提供更多的可靠性。然而，跟其他的东西一
样，随着分布式水平的提高，它们也会变得越来越复杂。系统发生错误的概率随着分布式水平的提
高呈指数级增长，因此即使简单的事情，如“如果断路器跳闸，则监控警报”，也就不那么简单
了。一个组件中的一个故障可能会在许多客户端和客户端的客户端上产生连锁反应，从而触发数千
个电路同时跳闸。而且，以前可能只需几行代码就能处理某个问题，而现在需要编写大量的代码才
能处理这些只存在于这个新世界的问题。 这也是大型复杂库，如Twitter的Finagle和Facebook的
Proxygen，深受欢迎的原因，它们能避免在每个服务中重写相同的逻辑。




                                                       4
大多数采用微服务架构的组织通常遵循上述模型。然而，随着系统中服务数量的增长，该模型逐渐
暴露出诸多问题。即便使用了诸如Finagle这样的库，项目团队仍需投入大量时间将这些库与系统
的其他部分集成，这一过程成本高昂且复杂。有时，这种成本显而易见，例如工程师被分配到专门
负责构建工具的团队；但更多时候，这种成本是隐性的，表现为产品研发周期的显著延长。
第二个问题在于，上述架构对微服务所使用的工具、运行时环境以及编程语言施加了限制。微服务
库通常针对特定平台开发，无论是某种编程语言还是像JVM这样的运行时环境。如果开发团队选
择了一个不受支持的平台，则往往需要将代码移植到新平台，这无疑是对有限工程资源的一种浪
费。在这种情况下，工程师无法专注于核心业务逻辑或产品功能，而是被迫投入大量精力构建工具
和基础设施。正因如此，许多中型企业倾向于在内部服务中仅支持单一平台，以减少不必要的复杂
性。
最后，值得一提的是该模型在管理层面的问题。尽管库模型可以抽象出实现微服务架构需求所需的
功能，但它本身仍然是一个需要维护的组件。必须确保数千个服务实例所使用的库版本一致或至少
兼容，而每次库的更新都会带来巨大的工作量：需要重新进行集成、测试和部署所有服务，即使这
些服务本身并未发生任何变化。这种管理负担不仅增加了运营复杂性，还可能导致服务更新的延
迟，进一步影响系统的敏捷性和可维护性。
正如我们在网络协议栈中所观察到的那样，大规模分布式服务所需的底层功能应当下沉至平台层实
现。通过这种方式，开发者能够基于高层协议（如HTTP）构建复杂的应用程序与服务，而无需关
注底层TCP协议如何管理网络数据包的传输细节。这种分层设计模式正是微服务体系结构所倡导
的：服务开发工程师可以将精力集中在业务逻辑的实现上，而不必耗费时间去重复编写服务基础设
施代码或管理整个系统的依赖库和框架。将这一理念映射到我们的架构图中，可以得到如下所示的
内容：




                                             5
通过修改网络协议栈来引入这一层次并非一项切实可行的任务。目前，许多解决方案倾向于利用一
组代理来实现该目标。其核心理念是，服务不会直接与下游组件建立连接，而是将所有流量导向一
个轻量级的软件组件，从而以透明的方式注入所需的功能。在这一领域，首个有据可查的技术突破
采用了“Sidecar”概念。Sidecar 是一个辅助进程，它与主应用程序并行运行，并为主应用提供额
外的功能支持。




随着微服务架构的日益普及，我们最近看到了一波新的代理浪潮，它们足以灵活地适应不同的基础
设施组件和偏好。
Service Mesh
在这种模式中，每个服务都配备了一个代理Sidecar 。由于这些服务只能通过代理Sidecar 进行通
信，我们最终会得到类似于下图的部署方案：




                                                  6
什么是服务网格（Service Mesh）？
服务网格是一种专用的基础设施层，旨在处理服务间通信。它通过复杂的服务拓扑结构，确保请求能够
以可靠且高效的方式进行传递。在实际应用中，服务网格通常以轻量级网络代理的形式实现，并以边车
（Sidecar）模式与应用程序代码共同部署。这种设计使得服务网格在运行过程中对应用程序完全透明，
无需修改业务逻辑或代码即可发挥作用。其核心价值在于突破了传统孤立代理的局限性，将分布式系统
中的代理组件视为一个有机整体，从而构建出一个具备可观测性、安全性和控制能力的统一网络层。这
种架构不仅提升了系统的可管理性，还为企业级微服务治理提供了强有力的支持。

随着微服务部署被迁移到更为复杂的运行时中去，如Kubernetes和Mesos，人们开始使用一些平
台上的工具来实现网格网络这一想法。他们实现的网络正从互相之间隔离的独立代理，转移到一个
合适的并且有点集中的控制面上来。




来看看这个鸟瞰图，实际的服务流量仍然直接从代理流向代理，但是控制面知道每个代理实例。控
制面使得代理能够实现诸如访问控制和度量收集这样的功能，但这需要它们之间进行合作：




                                                7
完全理解服务网格在更大规模系统中的影响还为时尚早，但这种架构已经凸显出两大优势。首先，
不必编写针对微服务架构的定制化软件，即可让许多小公司拥有以前只有大型企业才能拥有的功
能，从而创建出各种有趣的案例。第二，这种架构可以让我们最终实现使用最佳工具或语言进行工
作的梦想，并且不必担心每个平台的库和模式的可用性。
1.1.1.   终端节点
终端节点是网络的端点。它们包括个人计算机以及提供服务的服务器。终端节点专用于计算而非网
络，通常设计为支持通用计算。因此，终端节点瓶颈通常取决于结构（structure）和规模
（scale）。
 ● 结构：为了能够运行任意代码，个人计算机和大型服务器通常有一个操作系统在应用程序和硬
   件之间进行调解。为了简化软件开发，大多数大型操作系统被精心设计为分层软件（layered
   software）；为了保护操作系统免受其他应用程序的影响，操作系统实现了一套保护机制
   （protection mechanisms）；最后，核心操作系统例程（如调度程序和分配器）使用针对尽
   可能广泛的各类应用程序的通用机制（general mechanisms）。不幸的是，分层软件、保护
   机制和过度通用性的结合，即使使用最快的处理器，也会极大地减慢网络软件的速度。
 ● 规模：提供Web和其他服务的大型服务器的出现带来了进一步的性能问题。特别是，像Web服
   务器这样的大型服务器通常有成千上万的并发客户端。许多操作系统使用的数据结构和算法效
   率低下，这些算法是为连接数较少的时代设计的。
下表列出了主要终端节点瓶颈及其原因和解决方案。第一个瓶颈是由于传统操作系统结构导致数据
包数据在保护域之间进行拷贝；在Web服务器中，由于文件系统和校验和等其他操作（如检查所
有数据包数据的校验和）的类似拷贝，情况变得更加复杂。
                                                     8
 瓶颈             原因          示例解决方案
 拷贝             保护、结构       在没有操作系统干预的情况下拷贝多个数据
                            块（例如，RDMA）
 上下文切换          复杂调度        用户级协议实现 事件驱动的Web服务器
 系统调用           保护、结构       应用程序到驱动程序的直接通道（如 VIA）
 定时器            随定时器数量扩展    时间轮
 多路分解           随端点数量扩展     BPF和Pathfinder
 校验和/循环冗余校      通用性 随链路速度扩展 多比特计算
 验
 协议代码           通用性         头部预测

1.1.2.   中间设备
尽管我们专注于互联网路由器，但本书描述的所有技术同样适用于其他网络设备，如网桥、交换
机、网关、监视器和安全设备，以及IP之外的协议，如光纤通道。因此，在其余部分，将路由器视
为“通用网络互连设备”通常是有用的。与终端节点不同，这些是专用于网络处理的设备。因此，
路由器内部几乎没有结构开销，只使用非常轻量级的操作系统，并且有一个清晰分离的转发路径，
通常完全在硬件中实现。路由器面临的基本问题不是结构问题，而是由规模和服务引起的。
 ● 规模：网络设备面临两个方面的扩展：带宽扩展和节点数量扩展。带宽扩展是因为光链路不断
   变快。并且由于各种新应用的出现，互联网流量持续增长。节点数量扩展是因为随着更多企业
   上线，连接到互联网的端点越来越多。
 ● 服务：速度和规模的需求推动了20世纪80年代和90年代的网络行业，因为越来越多的企业上
   线（例如Amazon.com），并且创造了全新的在线服务（例如eBay）。但互联网的成功要求
   我们在未来十年更加关注通过提供性能、安全性和可靠性方面的保证来使其更有效—在拥塞时
   延迟、在攻击期间保护和在故障时可用性。
主要路由器（网桥/网关）瓶颈及其原因和解决方案。
 瓶颈       原因                  示例解决方案
 精确查找     链路速度扩展              并行哈希
 前缀查找     链路速度扩展 前缀数据库大小扩展 压缩多比特Trie树

                                                    9
 数据包分类    服务差异化 链路速度和大小扩展    决策树算法 硬件并行性（CAMs）
 交换       光电速度差距 排头阻塞        交叉开关 虚拟输出队列
 公平排队     服务差异化 链路速度扩展 内存扩   加权公平排队 赤字轮询 DiffServ，核心
          展                  无状态
 内部带宽     内部总线速度扩展           可靠条带化
 测量       链路速度扩展             Juniper的DCU
 安全       攻击数量和强度扩展          使用布隆过滤器的追踪 提取蠕虫签名

2.   硬件
如果链路速度为40Gbps（例如OC-768），路由器必须在8纳秒内完成40字节数据包的转发。在这种速
度下，数据包转发通常直接由硬件实现而非可编程处理器。若不了解硬件设计师的工具和约束条件，就
无法参与这类硬件密集型设计。不过，通过几个简单模型，即使对硬件毫无兴趣的人也能理解甚至参与
硬件设计。
互联网查找通常通过组合逻辑实现，数据包存储在路由器内存中，而路由器本身则由交换机和查找芯片
等组件构成。因此我们的旅程从逻辑实现开始，继而探讨内存内部结构，最后以基于组件的设计收尾。
更多细节可参考经典计算机体系结构著作（Hennessey和Patterson，1996）。
终端节点架构模型
处理器是一种状态机，它接收指令和数据序列作为输入，并将输出写入打印机和终端等I/O设备。为支持
大型程序，处理器将大部分状态存储在外部DRAM中。由于DRAM访问延迟较高（约60纳秒），处理器
通过SRAM缓存（小容量高速存储器）提升速度。一级缓存（L1）集成在处理器芯片上，二级缓存
（L2）位于外部。缓存通过地址哈希快速查找数据，利用时间局部性（频繁访问相同地址）和空间局部
性（访问相邻地址）提升效率。
如下图所示。处理器通过总线与其他组件通信，I/O设备采用内存映射机制（如网卡内存映射到地址100-
200），通过总线发送读写指令实现统一访问。现代机器支持直接内存访问（DMA），允许磁盘或网卡
直接读写内存，但总线竞争会导致处理器等待。图中，网卡实际位于与内存总线分离的I/O总线（如PCI
总线）上，其速度通常低于内存总线。




                                                       10
关键启示：网络应用的吞吐量受最慢总线（通常是I/O总线）限制。更糟的是，为保持操作系统结构，每
个收发包需要多次穿越总线。因此需要优化数据拷贝，避免冗余总线传输。
现代处理器采用多级流水线、超标量和多线程架构提升并行性（参见经典参考[Hennessey等，
1996]），但这些创新主要解决计算瓶颈而非数据移动瓶颈。下面通过示例探讨替代架构：
示例：基于交叉开关的终端节点架构




上图展示用可编程硬件开关（switch fabric）替代传统总线的设计。开关内部含多条并行总线，可同时
连接多对端点，以通过并行交换机实现并发处理和网络流量。例如处理器访问内存1时，网卡可同时访问
内存2。当处理器需读取网络数据时，开关可重新配置连接关系。若空包缓冲队列在两组内存间交替存
储，该设计能有效工作。Infiniband技术尝试替代处理器I/O总线。此例表明，即使协议设计者也能基于
简单硬件模型构想提升网络性能的架构。

                                                   11
路由器架构




一个路由器模型，标出了转发路径中的三个主要瓶颈：地址查找（B1）、交换（B2）和输出调度（B3）
模型涵盖高端（如Juniper M系列）和低端（如思科Catalyst）路由器。路由器本质是连接输入链路
（左）和输出链路（右）的盒子，其核心任务是根据包目的地址将其从输入链路交换到对应输出链路。
实际中两个方向的双向链路常集成在同一物理接口。路由器存在三大瓶颈：查找（B1）、交换（B2）和
输出调度（B3）。
查找
当包到达输入链路i时，处理器检查其32位IP地址。假设某包目的地址前六位是100100，处理器通过转发
表（FIB）确定输出链路。FIB存储前缀（如01*）与输出链路的映射关系，采用最长前缀匹配（LSP）规
则。例如100映射到链路6，1映射到链路2，则示例包应转发到链路6。




                                                  12
早期路由器采用共享处理器，后来（如思科GSR系列）改为每接口专用处理器。最初使用通用CPU，现
今最快路由器采用可编程ASIC芯片。但由于客户需求变化（如Web负载均衡），部分新路由器改用网络
处理器—专为网络优化的通用处理器。
交换
完成查找后，处理器控制内部交换系统将包从链路i转移到链路6。早期采用总线结构，但当N个输入链路
速率均为B时，总线带宽需达B·N。电气效应（如总线容性负载）限制了总线速度。现代高速路由器采用
并行交换架构，通过N条总线实现输入输出全连接。
队列管理
完成交换的包若遇链路拥塞，将进入输出队列。传统路由器采用FIFO队列，现代路由器通过复杂调度算
法保障公平性和延迟上限。
除核心功能外，路由器还需处理以下任务：
 ● 头部校验：检查版本号、选项长度，验证校验和，递减TTL并重算校验和。这些操作通常硬件实现。
 ● 路由处理：通过RIP/OSPF（域内）和BGP（域间）协议构建转发表。LSP包被识别为路由处理器目
    标后，会通过交换系统送达路由处理器。后者维护链路状态数据库并计算最短路径，最终将新转发
    表加载到各转发处理器。
 ● 协议处理：实现SNMP（提供远程可读计数器）、TCP/UDP（远程通信）和ICMP（错误报文）。
 ● 分片与ARP：部分功能如分片（大包拆分）、重定向（错误路由通知）和ARP（地址解析）可能仍
    在快路径处理。
现代路由器还需支持新功能：
 ● 内容感知处理：企业级路由器可能根据包内容（如Web URL）转发到不同服务器。
 ● 流量统计：计费和流量监测

示例：
可编程交换机
可编程交换机是一种网络设备，它具有能够根据用户需求进行自定义配置的功能。它通过使用软件定义
网络（SDN）技术，使用户能够灵活地控制和管理网络流量。与传统交换机相比，具有更高的灵活性和
可扩展性。传统交换机的功能是固定的，无法根据用户需求进行定制。而可编程交换机可以通过编程和
配置来实现各种不同的网络功能和策略。
📎
an_exhaustive_survey_on_p4_programmable_data_plane_switches_taxonomy_applications_challen
ges_and_future_tren.pdf
网络处理器（Network Processor）
网络处理器是可编程处理器，专为网络流量优化。例如英特尔IXP1200网络处理器（[Spalink等，2000]
评估）内含六个177MHz处理器，每个处理器从输入队列接收包，DRAM存储包数据，查找后将包标记输
                                                                                      13
出链路放入输出队列。华为也有大量的NP交换机。处理器需负责DRAM读写。IXP1200中，32字节数据
从队列到DRAM需45周期，反方向需55周期。40字节最小包处理需200周期=1.12微秒，仅支持约90万
包/秒的转发速率。通过六核并行和线程级并行（每个处理器维护四个上下文），理论峰值达21.6M包/
秒。部分网络处理器还提供特殊指令加速查表，或通过硬件引擎实现包头/数据分离处理（避免小包场景
下的性能瓶颈）。
缓冲与光交换
随着光纤链路速率提升，核心路由器中实现组合逻辑和存储的电子器件成为瓶颈。当前光信号需转换为
电信号进行路由处理，再转换回光信号发送。全光路由器可避免光电转换，但光学IP查找和密集光存储
仍不成熟。因此光通信创业公司多采用电子控制的光电路交换（非包交换），作为ISP网络核心连接传统
路由器。当R1-R2间流量增加时，运营商可人工调整光路带宽（时间尺度约分钟级）。但由于流持续时
间短导致资源浪费，包交换路由器短期内仍将占主导地位。

3.   协议
3.1.   抽象协议模型
协议是参与节点的状态机，包含接口和消息格式规范。下图展示了协议状态机模型，描述状态机如何响
应接口调用、接收消息和定时器事件（如发送消息、设置定时器）。




                   协议节点的状态机抽象模型
协议共有的功能：
 1. 底层数据操作：协议需收发数据包，涉及读写每个字节的操作（如TCP将接收数据复制到应用缓冲
    区，路由器将包从输入链路交换到输出链路）。TCP头还包含需计算数据校验和的操作，数据复制
    也涉及缓冲区等资源分配。
 2. 顶层解复用：状态机需将数据分发给多个客户端（如TCP根据端口号将网页数据分发给浏览器进
    程，并可能唤醒该进程）。


                                                    14
虽然这些通用功能代价高昂，但通过恰当技术可显著降低开销。衡量系统（如网络或单台路由器）性能
的两大核心指标是吞吐量（单位时间完成的作业数）和延迟（作业完成的最长时间）。ISP等系统所有者
追求最大吞吐量以增加收益，用户则希望端到端延迟低于几百毫秒（影响远程过程调用等计算速度）。
通信协议通常包含以下3个要素：
（1）语义：是指对构成协议的元素含义的解释。例如，在基本型数据链路控制协议中规定，数据
报文中的第一个协议元素的语义表示所传输报文的报头开始，接着是报文；而第二个协议元素的语
义表示正文的开始，接着是正文，等等。
（2）语法：规定将若干协议元素和数据结合起来表示一个完整内容时所应遵循的格式。如，在传
送数据报文时，协议元素和数据组合的次序为：报头开始符，报文，正文开始符，正文，正文结束
符，最后是奇偶校验码。
（3）规则：规定事件的执行顺序和匹配速率，如通信双方进行发送和应答的次序。作为典型的分
布式系统，网络系统通过协议保证整体系统工作。上层通过修改协议配置、策略和规则影响网
络。
3.2.   网络协议模型
为了保证通信双方能够正确，自动地进行通信，必须在信息传输顺序，信息格式和信息内容等方面制定
一套规则和约定，这种规则和约定的集合就是网络协议（network protocol）。
抽象导致分层。为支持不同通信介质，不同机器平台，多种通信规范，多种应用业务，要解决很多技术
问题。实现计算机网络通常采用层次结构，每一层完成特定功能，各层协调起来构成整个网络系统。总
的思想是从底层硬件提供的服务开始， 每一层提供一级服务，使得高层服务能够由低层服务来实现。换
句话说，每个协议需要定义两类接口：1）服务接口（service interface），定义本地对象可以在协议上
执行的操作，如发送接收消息操作；2）对等接口（peer interface），定义在实现通信服务的协议对等
实体间交换的信息格式/含义和通信方式。协议的定义通常我们需要去读 协议规范（protocol
specification）。协议规范用文字，伪码，状态转移图，分组格式化图等方式表示。程序员需要根据协
议规范实现可操作的（interoperate）程序。在网络体系结构中网络通信的建立必须是在通信双方的对等
层进行，不能交错。 在整个数据传输过程中，数据在发送端时经过各层时都要附加上相应层的协议头和
协议尾（仅数据链路层需要封装协议尾）部分，也就是要对数据进行协议封装，以标识对应层所用的通
信协议。




                                                     15
协议的设计过程中需要考虑几件事情。1）每一层功能定义；2）协议封装；3）跨层消息的复用和
解复用

开放互联模型（OSI）
OSI（Open System Interconnect），即开放式系统互联。 一般都叫OSI参考模型，是ISO（国际标准化
组织）组织在1985年研究的网络互联模型。该体系结构标准定义了网络互连的七层框架（物理层、数据
链路层、网络层、传输层、会话层、表示层和应用层），即ISO开放系统互连参考模型。在这一框架下进
一步详细规定了每一层的功能，以实现开放系统环境中的互连性、互操作性和应用的可移植性。
各层功能叙述如下：
(1)物理层(Physical Layer)
物理层是OSI参考模型的最低层，它利用传输介质为数据链路层提供物理连接。它主要关心的是通过物理
链路从一个节点向另一个节点传送比特流，物理链路可能是铜线、卫星、微波或其他的通讯媒介总的来
说物理层关心的是链路的机械、电气、功能和规程特性。
(2)数据链路层(Data Link Layer)
数据链路层是为网络层提供服务的，解决两个相邻结点之间的通信问题，传送的协议数据单元称为数据
帧。 数据帧中包含物理地址（又称MAC地址）、控制码、数据及校验码等信息。该层的主要作用是通过
校验、确认和反馈重发等手段，将不可靠的物理链路转换成对网络层来说无差错的数据链路。 此外，数
据链路层还要协调收发双方的数据传输速率，即进行流量控制，以防止接收方因来不及处理发送方来的
高速数据而导致缓冲器溢出及线路阻塞。
(3)网络层(Network Layer)

                                                            16
网络层是为传输层提供服务的，传送的协议数据单元称为数据包或分组。该层的主要作用是解决如何使
数据包通过各结点传送的问题，即通过路径选择算法（路由）将数据包送到目的地。另外，为避免通信
子网中出现过多的数据包而造成网络阻塞，需要对流入的数据包数量进行控制（拥塞控制）。当数据包
要跨越多个通信子网才能到达目的地时，还要解决网际互连的问题。
(4)传输层(Transport Layer)
传输层的作用是为上层协议提供端到端的可靠和透明的数据传输服务，包括处理差错控制和流量控制等
问题。该层向高层屏蔽了下层数据通信的细节，使高层用户看到的只是在两个传输实体间的一条主机到
主机的、可由用户控制和设定的、可靠的数据通路。 传输层传送的协议数据单元称为段或报文。
(5)会话层(Session Layer)
会话层主要功能是管理和协调不同主机上各种进程之间的通信（对话），即负责建立、管理和终止应用
程序之间的会话。会话层得名的原因是它很类似于两个实体间的会话概念。例如，一个交互的用户会话
以登录到计算机开始，以注销结束。
(6)表示层(Presentation Layer)
表示层处理流经结点的数据编码的表示方式问题，以保证一个系统应用层发出的信息可被另一系统
的应用层读出。如果必要，该层可提供一种标准表示形式，用于将计算机内部的多种数据表示格式
转换成网络通信中采用的标准表示形式。数据压缩和加密也是表示层可提供的转换功能之一。
(7)应用层(Application Layer)
应用层是OSI参考模型的最高层，是用户与网络的接口。该层通过应用程序来完成网络用户的应用
需求，如文件传输、收发电子邮件等。

OSI的设计者想当然的认为这个模型和以此模型为框架开发的协议将会主宰计算机通信领域，最终
取代各种专用的协议实现，并与像TCP/IP这样的多厂商模型竞争。只是这一设想从来都没实现
过。虽然在OSI背景下成功地开发了很多有用的协议，但是整体上的七层模型并未取得成功。相
反，TCP/IP体系结构已经占了主导地位。导致这一结果的原因有很多。也许其中最重要的原因在
于当关键性的TCP/IP协议已经成熟并经过了充分的测试的时候，类似的OSI协议却还在开发阶
段。就在企业开始认识到有必要进行跨网交互操作的时候，只有TCP/IP是现成的，并且已经准备
就绪。另一个原因是OSI模型没有必要这么复杂，它用七层结构实现的任务在TCP/IP中用了较少
的层照样完成。

TCP/IP协议体系结构
TCP／IP 协议体系结构是在实验性的分组交换网ARPANET上展开的协议研究开发工作的成果，
ARPANET由美国国防部高级研究计划署(Defense Advanced Research Projects Agency，DARPA)资
                                                                     17
助。TCP／IP协议体系结构通常称为TCP／IP协议族。这个协议族集合了大量的协议，这些协议已经通过
因特网体系结构研究会(Internet Architecture Board, IAB)作为因特网标准发布。TCP/IP网络体系结构，
并非国际标准，但由于其简洁，实用，在计算机网络中占有很重要的地位，成为事实上的工业标准。
TCP/IP的各层
大体上说，可以认为通信涉及三个要素：应用程序、计算机和网络。应用程序的例子包括文件传送程序
和电子邮件。在这里我们关心的应用程序是分布式的应用程序，涉及到两个计算机系统之间的数据交
换。这些应用程序以及其他一些应用程序，通常运行在能够支持多个并行程序的计算机上。计算机连接
到网络，需要交换的数据通过网络从一台计算机传送到另一台计算机上。因此，把数据从一个应用程序
传送到另一个应用程序，首先要做的就是设法把数据交给应用程序所在的计算机，还要设法将数据传送
给计算机中的目的应用程序。




头脑中有了这些概念，我们可以将通信任务分割成五个相对独立的层：
 ● 物理层：负责的是数据传输设备（如工作站或计算机）与传输媒体或网络之间的物理接口。这一层
   主要定义了传输媒体的特点、信号状态、数据率等诸如此类的特征。
 ● 链路层：负责封装和解封装IP报文，发送和接受ARP/RARP报文等。
 ● 网络层：负责路由以及把分组报文发送给目标网络或主机。网络接入层的主要任务是为与同一个网
   络相连的两个系统提供网络接入，并且为它们的数据选择路由以穿过网络。如果两个设备分别连接
   到不同的网络，这时需要有一些进程负责将它们的数据跨越多个互联的网络。这就是网际层的功
   能。这一层使用网际协议（Internet Protocol，IP)来提供经过多个网络的路由选择功能。IP协议不仅
   要在端系统中实现，同样也要在路由器中实观。路由器是连接两个网络的处理器，它的主要功能是
   把数据沿着从源到目的端系统的路径从一个网络转发到另一个网络。
 ● 传输层：负责对报文进行分组和重组，并以TCP或UDP协议格式封装报文。不管正在交换数据的应
   用程序是什么，通常都要求数据能够可靠地交换。也就是说，我们希望确保所有数据都能顺利到达
   目的应用程序，并且到达时的顺序与它们发送时的顺序是一致的。如同我们将要看到的，提供可靠
                                                                 18
    性的机制与应用程序的类型基本上没什么太大的关系。因此，有理由将这些机制集中到同一层中，
    并由所有的应用程序来共享。这一层就称为主机对主机层或运输层。提供此类功能的最常用的协议
    是传输控制协议(Transmission Control Protocol, TCP)。
  ● 应用层：负责向用户提供应用程序，比如HTTP、FTP、Telnet、DNS、SMTP等。应用层所包含的
    是用于支持各种不同的用户应用程序的逻辑。对于各种不同类型的应用程序，如文件传送程序，需
    要一个独立的专门负责该应用程序的模块。

4.   软件
4.1.   网络协议栈
操作系统是一种位于硬件之上的软件，旨在为应用程序员提供更便捷的开发环境。在大多数互联网路由
器中，对处理时间敏感的包转发直接在硬件上运行，而不经过操作系统。时间不敏感的代码则运行在精
简版的操作系统上，例如Cisco的IOS。然而，为了提高端到端性能（例如云计算），实现者需要理解操
作系统的成本与收益。

在操作系统设计中，宏内核（Monolithic Kernel）、微内核（Microkernel） 和 libOS（Library
Operating System） 是三种典型的架构模式。它们在系统结构、模块组织方式以及网络协议栈等子系
统的实现上存在显著差异。下面从基本概念出发，逐步分析它们在网络协议栈处理上的区别。

     基本概念
4.1.1.

1. 宏内核（Monolithic Kernel）
  ● 定义：将操作系统的核心功能（如进程管理、内存管理、文件系统、设备驱动、网络协议栈等）全
    部运行在内核态，作为一个整体运行。
  ● 特点：
     ○ 高性能：因为所有模块在内核空间直接调用，通信开销小。
     ○ 复杂性高：代码量大，模块耦合度高，稳定性较差，一个模块出错可能影响整个系统。
  ● 代表系统：Linux、Unix（如BSD）、Windows NT（虽然部分模块在用户态，但核心仍偏向宏内
    核）。
2. 微内核（Microkernel）
  ● 定义：只将最基本的功能（如进程调度、进程间通信、基本内存管理）放在内核态，其他功能（如
    文件系统、设备驱动、网络协议栈）放在用户态作为独立服务运行。

                                                               19
  ●  特点：
       ○ 模块化强：各功能模块独立，易于扩展和维护。
       ○ 稳定性高：一个服务崩溃不会导致整个系统崩溃。
       ○ 性能较低：因为模块间通信需要通过内核进行消息传递，开销较大。
  ● 代表系统：Mach、QNX、L4微内核系列。

3. libOS（Library Operating System）
  ● 定义：将操作系统功能以库的形式链接到应用程序中，每个应用拥有自己的“个性化”操作系统，运
     行在用户态，不依赖传统内核提供系统调用。
  ● 特点：
       ○ 高度定制化：每个应用可以按需选择需要的系统功能。
       ○ 性能潜力高：因为没有传统内核的系统调用开销，直接在用户态执行系统功能。
       ○ 隔离性差：每个应用独立运行自己的OS，缺乏统一的内核管理，隔离性和安全性需要额外机制
         保障。
  ● 代表系统：Google 的 gVisor、Unikernel（如 MirageOS）、Amazon 的 Firecracker（部分理
     念）。

     网络协议栈的处理区别
4.1.2.

网络协议栈是操作系统中的重要子系统，负责处理网络数据包的收发、协议解析、路由等任务。在不同
架构下，网络协议栈的实现方式、性能表现和系统交互方式存在显著差异。
1. 宏内核中的网络协议栈
  ● 位置：网络协议栈完全运行在内核态。
  ● 特点：
     ○ 高性能：因为协议栈与内核其他模块（如网络设备驱动）直接交互，无需跨特权级调用，数据
       路径短，延迟低。
     ○ 系统调用开销小：应用程序通过系统调用（如      read/ 
                                         write）访问网络，由内核直接处
       理。
     ○ 开发复杂：协议栈的修改和扩展需要重新编译内核或加载内核模块，灵活性较差。
  ● 典型实现：Linux 的 TCP/IP 协议栈就是典型的宏内核实现。

2. 微内核中的网络协议栈
  ● 位置：网络协议栈作为用户态的服务运行，与内核隔离。
  ● 特点：

                                                                      20
    ○    模块化与安全性高：协议栈崩溃不会影响内核，系统更稳定。
       ○ 性能较低：因为应用程序与协议栈之间需要通过内核的进程间通信（IPC）机制进行交互，数据
         需要在用户态和内核态之间多次拷贝，延迟较高。
       ○ 可扩展性好：可以动态加载或替换不同的网络协议栈实现。
  ● 典型实现：Mach 微内核中，网络协议栈是作为独立用户态服务运行的；QNX 也采用类似方式。

3. libOS 中的网络协议栈
  ● 位置：网络协议栈作为库函数直接链接到应用程序中，完全运行在用户态。
  ● 特点：
       ○ 极致性能潜力：没有系统调用和内核切换开销，数据路径最短，适合高性能场景（如高频网络
         通信、DPDK 类似场景）。
       ○ 完全定制化：应用可以自由选择或实现自己的协议栈，甚至可以针对特定硬件优化。
       ○ 隔离性与安全性挑战：由于缺乏内核级隔离，多个应用如果共享资源可能存在安全风险，需要
         依赖虚拟化或其他隔离机制（如沙箱）。
  ● 典型实现：
       ○ gVisor：虽然不是典型的 libOS，但它拦截系统调用并在用户态模拟内核行为，包括网络协议
         栈，从而提升安全性和可控性。
       ○ Unikernel：将整个应用和所需操作系统功能（包括协议栈）编译成一个单一镜像，运行在用户
         态或轻量虚拟化环境中，如 MirageOS。
       ○ DPDK + 应用程序：虽然不是严格意义上的 libOS，但 DPDK（Data Plane Development
         Kit）允许应用程序直接访问网卡，绕过内核协议栈，是一种“用户态协议栈”的实践。

抽象是我们为了应对现实世界的复杂性和不规则性而发明的理想化或假象。为了简化在裸机上的编程难
度，操作系统为应用程序员提供了抽象。处理原始硬件的三大核心难题是：中断处理、内存管理和I/O设
备控制。为了应对这些难题，操作系统提供了无中断计算、无限内存和简单I/O的抽象。
我们可以通过查看 /proc/interrupts 观察网卡的硬中断， /proc/softirqs 观察网卡软中断，
良好的抽象能提高程序员的生产力，但也会带来两个成本：
 1. 抽象机制的实现开销。例如，进程调度可能会增加Web服务器的开销。

 2. 抽象可能隐藏硬件能力，阻碍程序员优化资源使用。例如，操作系统的内存管理可能阻止程序员将
    查找数据结构保留在内存中以最大化性能。
接下来，我们将建模进程、虚拟内存和I/O 抽象的成本与底层机制。


                                                                 21
4.2.   通过进程实现无中断计算
程序在处理器上运行时，可能很快被网络适配器的中断打断。如果应用程序员必须直接处理中断，那么
一个100行的程序能正常运行都算奇迹。因此，操作系统提供了进程这一抽象，让程序员能够以无中断、
顺序执行的方式编写代码。
进程抽象通过三种机制实现：上下文切换、调度和保护，前两者下图所示。




程序员看到的进程P1似乎在独占处理器运行（上方时间线）。但实际上，处理器可能被定时中断打断，
导致操作系统调度程序接管（下方时间线）。此时，操作系统需保存P1的状态到内存，并可能切换到进
程P2运行，恢复其状态。因此，处理器的时间线可能涉及频繁的进程间上下文切换，由调度程序协调。
最后，保护机制确保一个进程的错误或恶意行为不会影响其他进程。
作为计算代理，“进程”有三种形式：
 1. 中断处理程序：用于处理紧急请求（如网络适配器收到消息），仅使用少量状态（如寄存器）。

 2. 线程：轻量级进程，共享相同内存（变量），切换成本低于完整进程。

 3. 用户进程：使用完整机器状态（内存+寄存器），进程切换开销较高。

以下示例展示这些概念如何影响终端节点的网络性能。
示例. BSD UNIX中的接收端活锁（Receiver Livelock）
在BSD UNIX中，数据包到达时触发中断，处理器保存当前进程（如Java程序）状态，跳转到中断处理程
序。该程序将数据包复制到内核的IP包队列，请求一个软件中断线程，然后退出。若无其他中断，调度
程序会优先运行该线程（而非用户进程）。




                                                 22
内核线程完成TCP/IP处理，将数据包放入目标应用程序的套接字队列。调度程序随后可能唤醒浏览器进
程。
在高负载下，系统可能陷入接收端活锁（Mogul & Ramakrishnan, 1997）：CPU将所有时间用于处理入
站包，却因应用程序无法运行而最终丢弃这些包。例如，连续的数据包到达会让最高优先级的中断处理
程序独占CPU，导致IP或套接字队列溢出。
此外，网络代码的延迟和吞吐量受“进程”激活时间影响。假设中断延迟约2μs（空调用），Linux进程上
下文切换约10μs。虽然看似微小，但千兆以太链路上，10μs内可到达30个最小包（40字节）。


4.3.   通过虚拟内存实现无限内存
在虚拟内存中，程序员看到的线性地址空间由编译器分配变量位置（如变量X位于虚拟地址1010）。该抽
象通过页表映射和按需调页实现。




                                                         23
如图 程序员看到的是连续的虚拟内存（上方），实际通过页表映射到物理内存和磁盘页（下方）。
虚拟地址需转换为物理地址。最简单的方法是偏移映射，但更灵活的方式是页表查询：将虚拟地址的高
位（如20位）作为页号，低位（12位）作为页内偏移。所有虚拟页可映射到任意物理页，未驻留的页
（如图中虚拟页2）标记为磁盘存储，访问时触发缺页异常，由OS从磁盘加载。
这种机制解决了存储分配和内存限制问题：OS只需管理固定大小的空闲页列表，程序员则可使用近乎无
限（仅受磁盘和地址位数限制）的内存空间。
地址转换可能拖慢指令执行（每次需额外访问页表）。现代CPU通过转译后备缓冲器（Translation
Look-aside Buffer，TLB）缓存最近使用的映射，由内存管理单元（MMU）硬件完成转换。页表还提
供进程间保护：非法访问会触发缺页异常，确保进程无法越权读写内存。
虽然路由器转发直接操作物理内存，但终端节点的网络代码依赖虚拟内存。虚拟内存可能带来TLB未命
中等开销，但也提供优化机会（如通过操作页表避免数据拷贝，详见1.3数据拷贝）。


4.4.   通过系统调用实现简单I/O
让应用程序员直接处理各类I/O设备的复杂性是不可行的。因此，操作系统将设备抽象为可读写的内存。




                                                     24
如图 程序员看到的磁盘、网络适配器等设备如同普通内存（上方），实际由内核通过设备驱动管理复杂
细节（下方）。
从简单I/O接口调用到实际设备读写（带完整参数）的代码称为设备驱动。若仅考虑抽象，驱动代码可放
在公共库中供应用调用。但像磁盘这样的共享设备必须受保护，错误进程可能破坏数据。因此，安全设
计要求仅内核（不可被应用破坏的受保护区域）能直接控制设备。
当浏览器需读取网页时，必须通过系统调用跨越应用-内核边界。系统调用是一种受保护的函数调用，硬
件指令会提升特权级别（内核模式），允许访问OS内部。调用返回后，应用恢复普通权限。系统调用比
函数调用更昂贵（涉及特权切换和参数检查），现代机器上可能耗时几微秒。

4.5.   Socket 接口
应用程序与操作系统的接口就是Socket接口，即 Socket 是操作系统提供的网络编程接口，用于实现进
程间的网络通信。它具有良好的可移植性，几乎所有现代操作系统（如 Linux、Windows、macOS 等）
都支持 Socket 接口，使得开发者可以编写跨平台的网络应用程序。




                                                    25
应用程序将可靠传输的任务委托给传输控制协议（TCP）。TCP的职责是为收发双方应用程序营造双向
共享数据队列的假象—事实上两者之间隔着不可靠的网络。发送方写入本地TCP发送队列的数据，应当
按相同顺序出现在接收方的TCP接收队列中，反之亦然。TCP通过将队列数据分割成段并重传未收到确
认（ACK）的段来实现这一机制。若应用程序（如视频会议）不需要可靠性保障，可选择用户数据报协
议（UDP）。UDP不保证可靠性，因此不需要ACK或重传。
TCP和UDP等传输协议通过互联网将数据段从发送节点传至接收节点，实际传输工作则委托给互联网路
由协议IP。每个IP包由TCP段加上包含目标IP地址的路由头组成。数据包从源到目的地的过程经由中间路
由器，路由表通过路由协议（尤其在拓扑变化时）动态构建。常用路由协议包括距离向量协议（如
RIP）、链路状态协议（如OSPF）和策略路由协议（如BGP）。




                                                 26
socket() 函数原型
 1   int socket(int domain, int type, int protocol);


该函数用于创建一个套接字描述符（socket descriptor），类似于文件描述符，后续的网络操作都通过
这个描述符进行。函数返回值为非负整数表示成功，-1 表示失败。
domain（协议族）
指定网络层使用的协议族，常用取值包括：
  ● AF_INET：IPv4 协议
  ● AF_INET6：IPv6 协议
  ● AF_UNIX 或 AF_LOCAL：Unix 域套接字（本地通信）
  ● AF_PACKET：底层数据包接口（用于原始套接字）

type（套接字类型）
指定传输层的通信语义，主要分为三种类型：
流套接字（Stream Socket）
  ● 类型值： SOCK_STREAM
  ● 协议：通常使用 TCP（传输控制协议）
  ● 特点：
     ○ 面向连接：通信前需建立连接（三次握手），通信结束后释放连接（四次挥手）
     ○ 可靠传输：保证数据顺序和完整性，提供丢包重传机制
     ○ 字节流：无消息边界，数据像水流一样连续传输
  ● 适用场景：HTTP、SMTP、FTP 等需要可靠传输的应用

数据报套接字（Datagram Socket）
  ● 类型值： SOCK_DGRAM
  ● 协议：通常使用 UDP（用户数据报协议）
  ● 特点：
     ○ 无连接：无需建立连接，直接发送数据
     ○ 不可靠：不保证数据到达顺序或是否到达
     ○ 数据报：有消息边界，接收方需按发送的完整数据报接收
     ○ 低延迟：无需握手和重传机制，传输速度快
  ● 适用场景：实时音视频流、游戏、DNS 等对延迟敏感的应用

原始套接字（Raw Socket）
                                                       27
 ●  类型值： SOCK_RAW
  ● 特点：
      ○ 直接访问底层协议（如 IP、ICMP）
      ○ 允许自定义协议头，可用于网络协议测试或实现特殊功能
      ○ 需要 root/admin 权限
  ● 适用场景：网络诊断工具（如 ping、traceroute）、协议分析器、网络攻击工具等

protocol（协议号）
网络协议栈根据端口号识别应用进程，将数据交给对应的应用程序。设置为 0，表示由系统根据 domai
n 和 type 自动选择合适的协议：

  ● 对于  SOCK_STREAM + AF_INET，默认使用 TCP（协议号 6）
  ● 对于  SOCK_DGRAM + 
                       AF_INET，默认使用 UDP（协议号 17）

若需指定特定协议（如 ICMP），可使用对应的协议号（如 IPPROTO_ICMP）。常用协议和端口号有




4.6 网络驱动模型
成功建立 socket 连接后，我们就能获得它的输入输出流，通信的本质是对输入输出流的处理。通
过输入流，读取网络连接上传来的数据，通过输出流，将本地的数据传出给远端。




                                                    28
在操作系统中，数据包的封装格式是 sk_buff（Socket Buffer）
在 Linux 内核网络子系统中， sk_buff是用于封装网络数据包的核心数据结构，负责在网络协议栈
（从网卡驱动到应用层）中传递数据包。它的设计直接影响网络性能和内存管理效率。
 sk_buff 定义于 
              <linux/skbuff.h>，用于存储数据包的 元数据（如协议头、长度、校验和
等）和 有效载荷（数据内容）。其核心字段包括：
关键字段




                                                           29
1    struct sk_buff {
2         用于管理数据包的长度
         /*
3          表示整个数据包的总长度（包括所有协议头和数据负载）。
         len:
4             ：表示当前协议层以下的数据长度（例如，在网络层，data_len 表示链路层头部的
         data_len
     长度）。
 5     通过 len 和 data_len 的配合，sk_buff 可以灵活地管理多层协议头和数据负载。
 6        */
 7        unsigned int len, data_len;
 8
 9                通过一组指针来管理数据包的实际内容，这些指针包括：
          //sk_buff
10        /*
11        head：指向缓冲区的起始地址（分配的内存起始位置）。
12        data：指向当前协议层数据的起始位置（例如，链路层头部之后的位置）。
13        tail：指向当前协议层数据的结束位置（数据负载的末尾）。
14        end：指向缓冲区的结束地址（分配的内存结束位置）。
15        */
16        unsigned char *head, *data, *tail, *end;
17
18       描述数据包的网络层、传输层等协议头信息 */
          /*
19                        // 数据包来自或发往的网络设备
          struct net_device *dev;
20        struct sock *sk;// 关联的套接字这个字段。
21                        // 主要用于传输层（如 TCP、UDP），用于将数据
     包与应用程序的套接字关联起来。
22        union {
23            struct tcphdr *th;
24            struct udphdr *uh;
25            struct iphdr *ip_hdr;
26            // 其他协议头指针
27        };
28
29        /* 时间戳、优先级、路由信息等         */
30        unsigned long timestamp;       // 记录数据包的时间戳，用于计算延迟、排序等。
31        __u32 priority;                // 表示数据包的优先级，用于 QoS（服务质量）管
     理。
32      struct dst_entry *dst;           // 路由信息，用于描述数据包的下一跳、出口设备
     等信息。
33
34      /* 用于协议栈各层之间的传递 */
35        struct sk_buff *next, *prev;   //   链表指针，用于组成队列
36
37   };




                                                                  30
为了便于组织数据结构本身， sk_buff组织成双向链表




内存布局
 ● 动态缓冲区： sk_buff 本身占用固定大小内存，数据部分通过       head 和 
                                                    data 指针指向动
   态分配的缓冲区。
 ● 头部预留空间： head 与    data 之间的区域用于协议栈各层（如链路层、网络层）添加协议头
   （例如 Ethernet 头、IP 头）。
 ● tail和end的空间是为了cacheline align对齐的问题导致的，这是操作系统的内容。

分配缓存变化的过程




                                                              31
1. 当 TCP 被要求传输某些数据时，它会按照特定标准分配缓冲区（如 TCP 最大段大小
   （mss）、对分散 / 聚集 I/O 的支持等）。
2. TCP 通过（skb_reserve）在缓冲区头部预留足够空间，以容纳所有层的报头（TCP、IP、链
   路层）。参数 MAX_TCP_HEADER 是所有层报头的总和，其计算考虑了最坏情况：由于
   TCP 层不知道传输将使用何种接口类型，因此为每一层预留了尽可能大的报头空间。它甚至
   考虑了多个 IP 报头的可能性（因为当内核编译支持 IP-over-IP 时，可能会有多个 IP 报
   头）。
3. TCP 有效载荷被复制到缓冲区中。请注意，图 2-8 只是一个示例。TCP 有效载荷可能以不同
   的方式组织；例如，它可以存储为片段。在第 21 章中，我们将了解分段缓冲区（通常也称为
   分页缓冲区）的样子。
4. TCP 层添加其报头。
5. TCP 层将缓冲区交给 IP 层，IP 层也会添加其报头。
6. IP 层将 IP 数据包交给相邻层，该层会添加链路层报头。



sk_buff 在网络协议栈中的流动
   sk_buff 是网络协议栈中数据包的核心载体，它贯穿整个协议栈，从数据包的接收、协议解析、路
由、转发，到最终的发送。以下是 sk_buff 在网络协议栈中的典型流动过程：
数据包接收（从网卡到协议栈）
 1. 网卡接收到数据包后，通过 DMA 将数据包写入内核的内存缓冲区。

                                                     32
2. 网卡驱动将数据包封装为一个 sk_buff 结构体，并设置 head、data、tail、end 等指针，指
   向数据包的起始和结束位置。
3. sk_buff 被传递到网络协议栈的链路层（如以太网驱动），链路层解析以太网头部，更新
   data 指针，跳过链路层头部。
4. 链路层将 sk_buff 传递到网络层（如 IPv4/IPv6），网络层解析 IP 头部，更新 data 指针，跳
   过网络层头部。
5. 网络层根据目标 IP 地址进行路由决策，决定是否转发或交给传输层处理。
6. 如果数据包是发给本机的，网络层将 sk_buff 传递到传输层（如 TCP/UDP），传输层解析传
   输层头部，更新 data 指针，跳过传输层头部。
7. 传输层将 sk_buff 交给应用程序（通过套接字）。



数据包发送（从应用程序到网卡）
 1. 应用程序通过系统调用（如 send/write）将数据发送到内核。
 2. 内核将数据封装为 sk_buff，设置 data 指针指向应用层数据。
 3. 传输层（如 TCP/UDP）添加传输层头部，更新 data 指针，跳过传输层头部。
 4. 网络层（如 IPv4/IPv6）添加网络层头部，更新 data 指针，跳过网络层头部。
 5. 链路层（如以太网）添加链路层头部和尾部（如 CRC 校验），更新 data 指针，跳过链路层
    头部。
 6. 网卡驱动将 sk_buff 中的数据通过 DMA 写入网卡，最终发送到网络中。



值得注意的是，在早期的 Linux 内核中，网络接口卡（NIC）接收数据包时完全依赖 硬件中断：
● 每当网卡收到一个数据包，就会触发一次中断，通知内核处理数据。
● 问题：在高流量场景下，频繁的中断会导致 CPU 负载过高（中断上下文切换频繁），且多个
  中断之间可能存在竞争，降低整体吞吐量。

NAPI（New API） 是 Linux 内核中用于优化网络数据包接收性能的一种机制，全称为 New API
for Network Interfaces。它诞生于 2.4 内核版本，旨在解决传统中断驱动接收数据包时的性能瓶
颈问题，尤其在高并发网络场景下表现更优。
NAPI 采用 中断驱动与轮询结合 的方式，核心逻辑如下：
 1. 中断触发：当网卡收到第一个数据包时，触发硬件中断，通知内核有数据到达。



                                                          33
2.  关闭中断并轮询：内核在中断处理程序中 临时关闭该网卡的中断，然后通过 轮询方式批量处
    理后续数据包，直到缓冲区数据处理完毕或达到预设的轮询次数上限。
 3. 重新开启中断：处理完批量数据后，重新开启网卡中断，等待下一次数据到达

通过这样做，可以
 ● 减少中断次数：通过批量处理数据包，显著降低高流量下的中断频率，减轻 CPU 负担。
 ● 均衡 CPU 资源：避免单个网卡中断过度占用 CPU，可结合内核调度机制将数据包处理分散
    到多个 CPU 核心。




NAPI 与硬件联合优化
NAPI 的高效运行依赖网卡支持 中断合并（Interrupt Coalescing） 或 RSS（Receive Side
Scaling，接收端扩展） 等特性：
 ● 中断合并：网卡在收到多个数据包后才触发一次中断，减少中断频率，即外设中断次数累计到
   一定阈值，或者到达某个timeout，再向CPU发送中断。
 ● RSS：将不同流的数据包分发到多个 CPU 核心，结合 NAPI 的多核轮询，进一步提升并行处
   理能力。

                                                            34
从数据拷贝的角度再看一遍数据传递的过程：
注意
 1. sk_buff本身不存储数据，通过指针指向真正的数据包内存空间
 2. 在各层传递时，只需要调整指针相应位置即可




在 Linux 网络子系统中，当网卡接收到一个网络数据包时，会通过 DMA（Direct Memory Access，直
接内存访问）将数据包写入内核预先分配的内存缓冲区，这些缓冲区组成一个环形队列，称为 rx_rin
g（接收环形缓冲区）。接下来，内核需要从 rx_ring 中取出数据包，并将其封装为
 sk_buff（socket buffer）结构体，以便交给网络协议栈进行进一步处理。

步骤 1：网卡将数据包写入 rx_ring
这一步是由网卡硬件完成的，属于数据接收的前置阶段：
 1. 网卡从网络中接收到一个数据包。

 2. 网卡通过 DMA 将数据包写入内核内存中的       rx_ring 缓冲区。

 3. 网卡更新  rx_ring 的硬件指针（如  head 指针），指向新写入的数据包位置。

此时，数据包已经存在于内核内存中，但内核尚未感知到它的存在。
步骤 2：内核检测到 rx_ring 中有新数据包

                                                           35
内核需要通过某种机制发现 rx_ring 中有新的数据包到达，这可以通过两种方式实现：
 1. 中断机制：
     ○ 网卡在写入数据包后，会触发一个硬件中断，通知内核有新数据包到达。
     ○ 内核的中断处理程序会响应中断，并开始处理             rx_ring 中的数据包。

 2. 轮询机制（NAPI，New API）：
     ○ 在高流量场景下，中断机制可能带来较大的开销。因此，Linux 引入了 NAPI 机制。
     ○ 在 NAPI 模式下，内核会定期轮询       rx_ring ，检查是否有新数据包到达，而不是依赖中断。

无论哪种方式，内核最终都会开始从 rx_ring 中读取数据包。
步骤 3：内核从 rx_ring 中取出数据包
内核通过 rx_ring 的软件指针（如 tail 指针）访问 rx_ring 中的缓冲区，取出数据包。具
体过程：
 1. 内核检查   rx_ring 的 tail 指针，确定下一个待处理的数据包位置。

 2. 内核从  rx_ring 的缓冲区中读取数据包的内容（包括链路层头部、网络层头部、传输层头部和
    数据负载）。
 3. 内核将  tail 指针向前移动，指向下一个待处理的缓冲区，以便后续的数据包可以被处理。

此时，数据包已经被内核“看到”，但还没有被封装为 sk_buff。
步骤 4：分配 sk_buff 结构体
在从 rx_ring 中取出数据包后，内核需要为这个数据包分配一个 sk_buff（socket buffer）结构
体。 sk_buff 是 Linux 网络子系统的核心数据结构，用于封装和管理网络数据包。
 1. 内核调用   alloc_skb() 函数（或类似的分配函数）从内存中分配一个         sk_buff 结构体。

 2. sk_buff 的分配通常是从内核的 SLAB 分配器（内存管理机制）中完成的，以确保高效的内存分
    配和回收。
此时， sk_buff 已经被创建，但还没有与数据包的内容关联。
步骤 5：将数据包内容拷贝或映射到 sk_buff
这一步是关键：内核需要将 rx_ring 中的数据包内容与刚刚分配的 sk_buff 关联起来。根据具
体的实现方式，这一步可能涉及以下两种方法之一：
方法 1：数据拷贝（Copy）
 ● 内核将   rx_ring 缓冲区中的数据包内容拷贝到        sk_buff 的缓冲区中。
 ● 这种方法简单直接，但会带来额外的 CPU 开销，因为需要将数据从一个内存区域拷贝到另一个内存
    区域。
 ● 在早期的 Linux 实现中，或者在某些特殊场景下，可能会使用这种拷贝方式。


                                                                  36
方法 2：数据映射（Zero-Copy 或 Direct Mapping）
 ● 更常见的现代实现中，内核会直接将            rx_ring 缓冲区的物理地址或虚拟地址映射到       sk_buff
    的缓冲区中，而不进行数据拷贝。
 ● 这种方式避免了数据拷贝的开销，提高了性能。
 ● 例如：
     ○ 在 
         NAPI + 
                 GRO（Generic Receive Offload）的实现中，内核可以直接将   rx_ring 中
       的数据包内容映射到 sk_buff。
     ○ 在高性能网络框架（如      AF_XDP 或  DPDK）中，数据包甚至可以直接从     rx_ring 传递
       到用户态，完全避免内核中的数据拷贝。
无论采用哪种方式，最终的结果是： sk_buff 的 data 指针指向了数据包的内容（包括链路层头
部、网络层头部等）。
步骤 6：设置 sk_buff 的元信息
在将数据包内容与 sk_buff 关联之后，内核需要设置 sk_buff 的各种元信息，以便协议栈能够
正确地处理数据包。这些元信息包括：
 1. 数据包长度：
     ○ 设置 sk_buff 的 
                      len 字段，表示整个数据包的总长度（包括链路层头部、网络层头部、
       传输层头部和数据负载）。
     ○ 设置 data_len 字段，表示当前协议层以下的数据长度（例如，在网络层，              data_len 表
       示链路层头部的长度）。
 2. 协议类型：
     ○ 设置 sk_buff 的 
                      protocol 字段，表示数据包的协议类型（如        ETH_P_IP 表示 IPv4
       数据包）。
 3. 网络设备信息：
     ○ 设置 sk_buff 的 
                      dev 字段，指向接收数据包的网络设备（如网卡）。

 4. 时间戳：
     ○ 设置 sk_buff 的 
                      timestamp 字段，记录数据包的接收时间，用于计算延迟、排序等。

 5. 其他信息：
     ○ 设置 sk_buff 的 
                      priority（优先级）、       dst（路由信息）等字段，为后续的协议栈
       处理提供支持。
步骤 7：将 sk_buff 传递给协议栈
完成 sk_buff 的封装后，内核会将 sk_buff 传递给网络协议栈进行进一步处理。具体过程如
下：
                                                                       37
1. 链路层处理：
    ○ 内核调用链路层（如以太网驱动）的处理函数，解析链路层头部（如 MAC 地址）。
    ○ 更新 
         sk_buff 的 
                    data 指针，跳过链路层头部。

2. 网络层处理：
    ○ 链路层将 sk_buff 传递给网络层（如 IPv4/IPv6），解析网络层头部（如 IP 地址）。
    ○ 更新 
         sk_buff 的 
                    data 指针，跳过网络层头部。
3. 传输层处理：
    ○ 网络层将 sk_buff 传递给传输层（如 TCP/UDP），解析传输层头部（如端口号）。
    ○ 更新 
         sk_buff 的 
                    data 指针，跳过传输层头部。

4. 应用层处理：
    ○ 如果数据包是发给本机的，传输层会将       sk_buff 交给对应的套接字（ 
                                                  struct sock），
      最终传递给应用程序。

从 rx_ring 到 sk_buff 的关键点
 步骤 操作                        说明
 1      网卡将数据包写入 rx_ring    硬件完成，通过 DMA 写入内核内存
 2      内核检测到 rx_ring 中有新数据 通过中断或轮询机制
        包
 3      内核从 rx_ring 中取出数据包 通过 tail 指针访问缓冲区
 4      内核分配 sk_buff 结构体    从内存中分配一个新的 sk_buff
 5      将数据包内容映射或拷贝到 sk_buf 避免拷贝（零拷贝）或直接拷贝
      f

6     设置 sk_buff 的元信息        包括长度、协议类型、时间戳等
7     将 sk_buff 传递给协议栈       从链路层到应用层的逐层处理




                                                              38
